\let\negmedspace\undefined
\let\negthickspace\undefined
\documentclass[journal,12pt,onecolumn]{IEEEtran}
\usepackage{cite}
\usepackage{amsmath,amssymb,amsfonts,amsthm}
\usepackage{algorithmic}
\usepackage{graphicx}
\graphicspath{{./figs/}}
\usepackage{textcomp}
\usepackage{xcolor}
\usepackage{txfonts}
\usepackage{listings}
\usepackage{enumitem}
\usepackage{mathtools}
\usepackage{gensymb}
\usepackage{comment}
\usepackage{caption}
\usepackage[breaklinks=true]{hyperref}
\usepackage{tkz-euclide} 
\usepackage{listings}
\usepackage{gvv}                                        
%\def\inputGnumericTable{}                                 
\usepackage[latin1]{inputenc}     
\usepackage{xparse}
\usepackage{color}                                            
\usepackage{array}                                            
\usepackage{longtable}                                       
\usepackage{calc}                                             
\usepackage{multirow}
\usepackage{multicol}
\usepackage{hhline}                                           
\usepackage{ifthen}                                           
\usepackage{lscape}
\usepackage{tabularx}
\usepackage{array}
\usepackage{float}
%\newtheorem{theorem}{Theorem}[section]
%\newtheorem{theorem}{Theorem}[section]
%\newtheorem{problem}{Problem}
%\newtheorem{proposition}{Proposition}[section]
%\newtheorem{lemma}{Lemma}[section]
%\newtheorem{corollary}[theorem]{Corollary}
%\newtheorem{example}{Example}[section]
%\newtheorem{definition}[problem]{Definition}
\begin{document}

%\textbf{\Large 12.319} \\
%\textbf{\large AI252BTECH11033} \\
\title{12.319}
\author{AI25BTECH11033 - Spoorthi N}
% \maketitle
% \newpage
% \bigskip
%\begin{document}
{\let\newpage\relax\maketitle}
%\renewcommand{\thefigure}{\theenumi}
%\renewcommand{\thetable}{\theenumi}

\textbf{Question}:

Let $M$ be the real vector space of $2 \times 3$ matrices with real entries.  
Let $T : M \to M$ be defined by
\begin{align}
T\!\left( 
\begin{bmatrix}
x_{1} & x_{2} & x_{3} \\
x_{4} & x_{5} & x_{6}
\end{bmatrix}
\right)
=
\begin{pmatrix}
-x_{6} & x_{4} & x_{1} \\
x_{3} & x_{5} & x_{2}
\end{pmatrix}.
\end{align}

The determinant of $T$ is \_\_\_\_ \hfill (MA 2013)

\textbf{Solution}:

We are given the transformation
\begin{align}
T\!\left(
\begin{pmatrix}
x & y & z \\
u & v & w
\end{pmatrix}
\right)
=
\begin{pmatrix}
-w & u & x \\
z & v & y
\end{pmatrix}.
\end{align}

Represent input as a vector;

\begin{align}
\begin{pmatrix}
x & y & z \\
u & v & w
\end{pmatrix}
\quad \longleftrightarrow \quad
\begin{bmatrix}
x \\ y \\ z \\ u \\ v \\ w
\end{bmatrix}.
\end{align}

Transformation on the vector;

\begin{align}
T\!\left(
\begin{pmatrix}
x \\ y \\ z \\ u \\ v \\ w
\end{pmatrix}
\right)
=
\begin{pmatrix}
-w \\ u \\ x \\ z \\ v \\ y
\end{pmatrix}.
\end{align}

Matrix of transformation;

\begin{align}
A =
\begin{pmatrix}
0 & 0 & 0 & 0 & 0 & -1 \\
0 & 0 & 0 & 1 & 0 & 0 \\
1 & 0 & 0 & 0 & 0 & 0 \\
0 & 0 & 1 & 0 & 0 & 0 \\
0 & 0 & 0 & 0 & 1 & 0 \\
0 & 1 & 0 & 0 & 0 & 0
\end{pmatrix}.
\end{align}

Determinant:

\begin{align}
\det(A) = (-1)\times(-1) = 1.
\end{align}

\begin{align}
{\det(T) = 1}
\end{align}

\end{document}
